%% annex/p02/propositional-metatheory.tex --- proofs for P2.1 metatheory

\begin{annexitem}{annex:p2.3.1}{thm:p02-nd-soundness}{Soundness of classical propositional natural deduction}
\begin{proof}
Induct on derivations, strengthening the claim to every subderivation: a
valuation satisfying its open assumptions satisfies its conclusion.
Assumption, the conjunction rules, disjunction introduction, implication and
negation elimination, and bottom elimination follow from their truth clauses.
For disjunction elimination, a true major premise selects a true disjunct, and
the corresponding induction hypothesis gives the common conclusion.
For implication introduction, a true antecedent gives the consequent by
induction; a false one makes the implication true.  Negation introduction
would otherwise make \(\bottom\) true.  Under RAA, a false conclusion makes
its discharged negation true, with the same impossibility.  Hence every rule
preserves the strengthened claim.
\end{proof}
\end{annexitem}

\begin{annexitem}{annex:p2.3.2}{thm:p02-lk-soundness}{Soundness of the classical propositional sequent calculus}
\begin{proof}
Induct on the sequent derivation.  An identity sequent is satisfied whichever
value its common formula receives, and a \(\bottom\)-initial sequent has a
false antecedent member.  Suppose a valuation falsifies the conclusion of a
rule: every antecedent member is then true and every succedent member false.
Weakening therefore leaves the corresponding premise falsified; contraction
does not change supports.  For Cut, satisfaction of the left premise forces
the cut formula true, while satisfaction of the right premise forces it
false, a contradiction.

For each logical rule the truth clause for its principal connective falsifies
a premise.  Negation moves a false succedent formula to a true negated
antecedent, or a true antecedent formula to a false negated succedent.
For conjunction-left both conjuncts are true; for conjunction-right a false
conjunction makes one of the two premises false.  A true disjunction on the
left makes one disjunct true, while a false disjunction on the right makes
both disjuncts false.  A true implication on the left has a false antecedent
or true consequent, falsifying the corresponding premise; a false implication
on the right has a true antecedent and false consequent.  Hence no rule can
have universally satisfied premises and a falsifiable conclusion.
\end{proof}
\end{annexitem}

\begin{annexitem}{annex:p2.3.3}{thm:p02-mcs-closure}{Closure properties of maximal consistent sets}
\begin{proof}
If \(\Delta\derives\varphi\) but \(\varphi\notin\Delta\), maximality makes
\(\Delta\cup\{\varphi\}\) inconsistent.  Substituting the derivation for
each added assumption derives \(\bottom\) from \(\Delta\), proving clause 1.

Consistency forbids both \(\varphi\) and \(\neg\varphi\).  If the former is
absent, maximality derives \(\bottom\) after adjoining it.  Some open
occurrence remains; merge all under a fresh label and apply negation
introduction.  Clause 1 yields \(\neg\varphi\in\Delta\), without vacuous
discharge.  This proves clause 2.

The conjunction rules give clause 3.  Disjunction introduction gives one
direction of clause 4; for the other, absence of both disjuncts supplies both
negations by clause 2, and disjunction elimination derives \(\bottom\).

For clause 5, implication elimination handles the forward direction.  If
\(\psi\in\Delta\), assume \(\varphi\), conjoin it with \(\psi\), eliminate
back to \(\psi\), and discharge the genuinely used assumption.  If
\(\varphi\notin\Delta\), clause 2 supplies its negation; the assumed
\(\varphi\) yields \(\bottom\), then \(\psi\), and is discharged.  Clause 1
places either implication in \(\Delta\).
\end{proof}
\end{annexitem}

\begin{annexitem}{annex:p2.3.4}{thm:p02-lindenbaum}{Countable propositional Lindenbaum lemma}
\begin{proof}
From \(\Gamma\), enumerate formulas and adjoin each iff consistency remains;
let \(\Delta\) be their union.  Any \(\bottom\)-derivation uses finitely many
assumptions, hence lies in one consistent stage, impossible.  An omitted
formula was rejected, so adjoining it to \(\Delta\) is inconsistent.  Thus
\(\Delta\) is maximal consistent.
\end{proof}
\end{annexitem}

\begin{annexitem}{annex:p2.3.5}{thm:p02-truth-lemma}{Propositional truth lemma}
\begin{proof}
Induct structurally.  Variables follow by definition and \(\bottom\) by
consistency.  Negation uses clause 2 of Lemma~\ref{thm:p02-mcs-closure};
conjunction, disjunction, and implication use clauses 3--5 with their truth
clauses.  These exhaust the grammar.
\end{proof}
\end{annexitem}

\begin{annexitem}{annex:p2.3.6}{thm:p02-nd-completeness}{Strong completeness of classical propositional natural deduction}
\begin{proof}
Contrapositively, assume
\(\Gamma\notderives_{\mathsf{ND}_{\mathrm{CPL}}}\varphi\).  Then
\(\Gamma\cup\{\neg\varphi\}\) is consistent: a \(\bottom\)-derivation with
no open \(\neg\varphi\) yields \(\varphi\) by bottom elimination; otherwise
merge every such occurrence under a fresh label and apply RAA.  Either case
contradicts the assumption while respecting nonvacuous discharge.  Extend to
a maximal \(\Delta\).  Its truth valuation satisfies \(\Gamma\) and
\(\neg\varphi\), hence falsifies \(\varphi\).  Thus
\(\Gamma\not\entails\varphi\).
\end{proof}
\end{annexitem}

\begin{annexitem}{annex:p2.3.7}{thm:p02-compactness}{Propositional compactness, fixed countable language}
\begin{proof}
Completeness yields a finite derivation; its finite open-assumption set
\(\Gamma_0\) entails \(\varphi\) by soundness.  Applying this to
\(\Gamma\entails\bottom\) proves the contrapositive of finite
satisfiability.  Satisfiability passes to every subset.
\end{proof}
\end{annexitem}

\begin{annexitem}{annex:p2.3.8}{thm:p02-lk-completeness}{Completeness of the propositional sequent calculus}
\sketchnote{the rule-by-rule translation trees from natural deduction to
the registered finite-multiset LK calculus are omitted; the sketch uses
primitive Cut and does not establish cut elimination.}
Soundness is Theorem~\ref{thm:p02-lk-soundness}.  For the converse, induction
on a natural-deduction derivation translates introductions to corresponding
right rules, eliminations to left rules followed when necessary by Cut, and
discharge to right rules; weakening and contraction align finite contexts.
Thus \(\Gamma_0\derives_{\mathsf{ND}_{\mathrm{CPL}}}\chi\), for finite
\(\Gamma_0\), yields \(\Gamma_0\Rightarrow\chi\) in LK.

Assume the displayed semantic condition for \(\Gamma\Rightarrow\Delta\),
and list \(\operatorname{supp}(\Delta)=\{\delta_1,\ldots,\delta_m\}\),
allowing \(m=0\).  The set consisting of
\(\operatorname{supp}(\Gamma)\) and
\(\neg\delta_1,\ldots,\neg\delta_m\) is unsatisfiable.  Natural-deduction
completeness derives \(\bottom\) from finitely many of these formulas, and
the translation subroutine gives the corresponding LK sequent.  Cut with the
registered \(\bottom\)-initial sequent empties its succedent.  For each used
\(\neg\delta_j\), the identity and negation-right rows derive a sequent with
\(\delta_j,\neg\delta_j\) in the succedent; Cut moves \(\delta_j\) to the
target succedent.  Weakening and contraction then give exactly
\(\Gamma\Rightarrow\Delta\).  When \(m=0\) no such moves occur; empty
antecedents are covered as well.  No choice principle enters the countable
Lindenbaum construction.
\end{annexitem}

\begin{annexitem}{annex:p2.3.9}{thm:p02-interpolant-truth-table-correctness}{Correctness of the truth-table interpolant construction}
\begin{proof}
Let \(\theta\) be the output and let \(X\) be the common-variable set.  If a
valuation \(v\) satisfies \(\varphi\), its restriction \(v|_X\) is among the
selected assignments, so \(v\) satisfies the corresponding complete
conjunction and hence \(\theta\).  Thus
\(\varphi\entails\theta\).

Suppose instead that some valuation \(v\) satisfies \(\theta\) and falsifies
\(\psi\).  It satisfies a selected conjunction, whose assignment
\(u:X\to\{0,1\}\) therefore extends to a valuation \(w\) satisfying
\(\varphi\).  Amalgamate \(w\) on the variables occurring only in
\(\varphi\), \(v\) on those occurring only in \(\psi\), and their common
value \(u\) on \(X\).  The resulting valuation satisfies \(\varphi\) and
falsifies \(\psi\), contrary to the hypothesis.  Hence
\(\theta\entails\psi\).  By construction, every variable of \(\theta\) lies
in \(X\).  The empty selected family gives \(\bottom\), and an empty
conjunction gives \(\topv\), so the boundary cases obey the same argument.
\end{proof}
\end{annexitem}

\begin{annexitem}{annex:p2.3.10}{thm:p02-resolution-soundness}{Soundness of propositional resolution}
\begin{proof}
A valuation falsifying the resolvent falsifies every nonpivot literal.  To
satisfy both parents it must satisfy \(L\) and \(\overline L\), impossible.
Hence each resolution step preserves satisfaction; induction proves the claim
because every valuation falsifies the empty clause.
\end{proof}
\end{annexitem}

\begin{annexitem}{annex:p2.3.11}{thm:p02-resolution-completeness}{Refutational completeness of propositional resolution}
\begin{proof}
Soundness gives the reverse implication.  First suppose \(S\) is finite,
delete tautologies, and induct on the number of occurring variables.  With
none, unsatisfiability forces the empty clause.  Otherwise choose an occurring
\(p\); form \(S_p\) from the \(p\)-free clauses and all non-tautological
resolvents on \(p\).

Soundness gives one satisfiability direction.  Conversely, let \(w\) satisfy
\(S_p\).  A positive-pivot parent with false remainder requires \(p=1\); a
negative one requires \(p=0\).  Both cannot occur: their resolvent would be
false under \(w\), hence non-tautological, yet belongs to \(S_p\).  Choose a
permitted value; the extension satisfies \(S\).  Thus satisfiability is
preserved.  Since \(S_p\) has fewer variables, induction refutes it when
\(S\) is unsatisfiable; splice the one-step derivations of added resolvents to
obtain a refutation from \(S\).  This is Davis--Putnam elimination.

For arbitrary unsatisfiable \(S\), apply fixed-language compactness to the
formula interpretations of its clauses.  Choose one preimage in \(S\) for
each member of the finite unsatisfiable image subset; the resulting finite
\(S_0\subseteq S\) is unsatisfiable although interpretation need not be
injective.  The finite argument refutes \(S_0\), hence \(S\).  Compactness,
not the elimination induction, removes finiteness.
\end{proof}
\end{annexitem}

\begin{annexitem}{annex:p2.3.12}{thm:p02-dual-horn-sat}{Dual-Horn-SAT tractability}
\begin{proof}
Complement every variable throughout the instance.  A clause with at most one
negative literal becomes a clause with at most one positive literal, and the
map is a linear-time bijection between satisfying assignments (complement the
assignment in the same way).  Apply the linear-time Horn algorithm and map its
answer back.
\end{proof}
\end{annexitem}
